\documentclass[conference]{IEEEtran}
\usepackage[utf8]{inputenc}
\usepackage[T1]{fontenc}
\usepackage[spanish,es-nodecimaldot]{babel}
\usepackage{cite}
\usepackage{amsmath,amssymb}
\usepackage{graphicx}
\usepackage{booktabs}
\usepackage{xcolor}
\usepackage{url}
\usepackage{hyperref}
\hypersetup{colorlinks=true,urlcolor=blue,citecolor=blue,linkcolor=blue}

\title{Alertas de empleo remoto para Colombia con n8n: prototipo y evaluación del filtro geográfico}

\author{\IEEEauthorblockN{Andrés Daniel Medina Erazo}
\IEEEauthorblockA{Universidad del Cauca\\
Popayán, Colombia}}

\begin{document}
\maketitle

\begin{abstract}
Se presenta un prototipo en n8n que busca facilitar la búsqueda de empleo desde Colombia mediante un correo diario con hasta cinco vacantes remotas de LinkedIn. El flujo requiere evidencia de elegibilidad geográfica para incluir una oferta y deriva los casos ambiguos a revisión. Se evaluó este criterio en una muestra dirigida de 16 vacantes públicas de Remotavia consultadas el 24 de noviembre de 2026. Ocho admitían Colombia o tenían apertura mundial y ocho estaban restringidas a otros países o regiones. La regla de referencia aceptó 16 vacantes, de las cuales ocho eran elegibles. La regla basada en etiquetas explícitas seleccionó las ocho elegibles. La precisión pasó de 50\% a 100\% y la sensibilidad fue de 100\% con ambas reglas en la muestra. Quedan por evaluar la clasificación de anuncios de LinkedIn y el ahorro de tiempo para los usuarios del prototipo.
\end{abstract}

\begin{IEEEkeywords}
n8n, automatización, vacantes remotas, elegibilidad geográfica, Colombia
\end{IEEEkeywords}

\section{Introducción}
El taller propone automatizar parte de la búsqueda de empleo desde Colombia con un flujo de n8n que consulta ofertas de LinkedIn con Apify, limpia campos y prepara un correo diario con hasta cinco resultados y enlaces [1]. Su filtro geográfico acepta ubicaciones que contienen ``Colombia'', ``Remote'', ``Bogot'', ``Medell'' o ``Remoto''. La condición ``Remote'' expresa modalidad de trabajo. El anuncio puede limitar la contratación a países específicos.

El objetivo del prototipo es reunir oportunidades pertinentes en un solo correo. La evaluación aborda una pregunta concreta: ¿cuántas ofertas elegibles conserva el resumen y cuántas restringidas evita cuando se exige evidencia geográfica? Se entrega un flujo de alertas con clasificación en tres estados y una comparación reproducible sobre anuncios públicos.

\section{Método}
Se tomaron 16 vacantes del índice de empleo remoto de Remotavia, que muestra etiquetas de los países admitidos [2], [3]. La captura se realizó el 24 de noviembre de 2026. La selección fue dirigida y equilibrada: ocho anuncios etiquetados ``Colombia only'' u ``Open worldwide'' y ocho restringidos a Estados Unidos, Canadá, México, Brasil o regiones que excluían a Colombia. Se registraron título, empresa, texto de ubicación, etiqueta, clase de referencia y URL en el archivo CSV del proyecto. La etiqueta del índice se usó como referencia de elegibilidad.

Se compararon dos reglas. La regla base aceptó todo anuncio cuya ubicación incluyera ``Remote'' o ``Remoto'', equivalente al criterio amplio del flujo recibido. La regla mejorada aceptó solo las etiquetas ``Colombia only'' y ``Open worldwide''. Para cada anuncio se calcularon verdaderos positivos (TP), falsos positivos (FP), falsos negativos (FN) y verdaderos negativos (TN). Las métricas se definieron como precisión $TP/(TP+FP)$, sensibilidad $TP/(TP+FN)$ y exactitud $(TP+TN)/n$.

El prototipo desconecta la evaluación de Apify y Gmail para que la comparación pueda ejecutarse sin credenciales. La adaptación operativa conserva la búsqueda de LinkedIn, clasifica evidencia explícita en tres estados y deriva a revisión los casos ambiguos.

\section{Resultados}
Las 16 ubicaciones contenían la palabra ``Remote''. La regla base, por tanto, aceptó los 16 anuncios. Ocho eran elegibles para Colombia y ocho no lo eran. La regla base tuvo ocho falsos positivos. La regla mejorada aceptó los ocho elegibles, rechazó los otros ocho y no tuvo errores en esta muestra.

\begin{table}[htbp]
\caption{Clasificación de las reglas sobre 16 vacantes}
\label{tab:resultados}
\centering
\begin{tabular}{lrrrrrr}
\toprule
Regla & TP & FP & FN & TN & Prec. & Exact. \\
\midrule
Remote/Remoto & 8 & 8 & 0 & 0 & 50\% & 50\% \\
Elegibilidad explícita & 8 & 0 & 0 & 8 & 100\% & 100\% \\
\bottomrule
\end{tabular}
\end{table}

La sensibilidad fue 100\% para ambas reglas. En este benchmark, requerir una etiqueta de elegibilidad elevó la precisión de 50\% a 100\% y redujo los falsos positivos de ocho a cero. No se interpreta este cambio como una estimación de desempeño en el conjunto completo de LinkedIn.

\section{Discusión}
La comparación aísla una diferencia entre modalidad y lugar de contratación. Su resultado depende de que el rótulo de elegibilidad refleje correctamente la vacante. La muestra es pequeña, seleccionada para contener casos positivos y negativos, y proviene de un agregador que normaliza anuncios de varias fuentes. Por ello, no permite estimar prevalencia ni generalizar porcentajes al mercado laboral.

El flujo mejorado distingue evidencia positiva explícita, restricciones a otros países y casos inciertos. Un anuncio de LinkedIn puede no exponer una etiqueta normalizada equivalente. En ese caso, la clasificación textual debe conservar ``revisar'' si no hay evidencia suficiente; un resultado automatizado no confirma por sí mismo condiciones contractuales. La conexión real a Apify requiere credenciales externas y no se ejecutó para esta evaluación.

\section{Conclusión}
El prototipo reúne en un correo diario enlaces a vacantes que muestran evidencia de admitir postulaciones desde Colombia. En la muestra de Remotavia, exigir una etiqueta explícita dejó fuera ocho ofertas restringidas y conservó las ocho vacantes elegibles. Se entrega un flujo n8n importable para reproducir la comparación y una adaptación operativa para preparar el correo. El siguiente paso es contrastar la clasificación del flujo con una revisión manual de anuncios de LinkedIn y medir su efecto en el tiempo de búsqueda.

\begin{thebibliography}{00}
\bibitem{guia} J. A. Zúñiga Daza, B. H. Meneses Macías y E. L. Restrepo Gaviria, ``Taller n8n: Alerta Empleo Colombia LinkedIn. Guía paso a paso detallada,'' material de taller suministrado, s. f.
\bibitem{co} Remotavia, ``Remote jobs in Colombia,'' consultado el 24 de noviembre de 2026. [En línea]. Disponible: \url{https://remotavia.com/remote-jobs-colombia}
\bibitem{all} Remotavia, ``Remote jobs, checked for who can actually apply,'' consultado el 24 de noviembre de 2026. [En línea]. Disponible: \url{https://remotavia.com/remote-jobs}
\bibitem{data} A. D. Medina Erazo, ``Benchmark de elegibilidad geográfica para vacantes remotas,'' conjunto de datos y código del proyecto, captura del 24 de noviembre de 2026. [En línea]. Disponible: \url{https://github.com/DANIELMEDINA0408/filtro-elegibilidad-vacantes-colombia}
\bibitem{manual} Universidad del Cauca, ``Manual de Identidad Visual,'' versión julio de 2024. [En línea]. Disponible: \url{https://www.unicauca.edu.co/wp-content/uploads/2024/08/Manual-de-Identidad-Visual_Unicauca_V2024-1.pdf}
\end{thebibliography}

\end{document}
